\section{Related Work}
\label{sec:related_work}

The related works presented in this section discuss studies connected to the context of quality indicators in web applications based on Core Web Vitals (CWV) metrics, evaluation and comparison of performance aspects, and identification and comparison of computational costs in web applications.

\citet{parsingJson:2019} present  \textit{simdjson} tool, a high-performance solution for processing large volumes of data in JSON. The study emphasizes that, with the growing exchange of structured data in web applications, the performance of traditional parsers has become a bottleneck. \textit{Simdjson} outperforms reference analyzers such as RapidJSON by employing optimized techniques for detecting JSON structures. Thus, the study conceptually contributes to the present study by highlighting how the interpretation and manipulation of JSON data can directly impact overall application performance, particularly in contexts where different rendering strategies are compared.

\citet{webPerformanceOptimize:2023} analyzed an industrial case study in an agricultural technology company aiming to improve the performance of a large-scale web dashboard. The authors developed a performance engineering plan using a benchmarking tool based on Google Lighthouse to assess the impact on web performance metrics. The results demonstrated a significantly faster browsing experience, with pages loading 98.37\% faster on desktop devices and 97.56\% faster on mobile. Content rendering also became more efficient, with improvements of 48.25\% and 19.85\% in rendering speed on desktop and mobile devices, respectively. This study is relevant to the current research as it provides data on the impact of different interventions affecting performance metrics using Google Lighthouse. However, there remains a gap for analyzing performance during data transactions between JSON and HTML formats using the same tool.

\citet{webAutoscaling:2022} explored mechanisms to ensure the resilience and performance of web applications under high-demand scenarios and potential failures. The study investigated different scaling approaches for applications and databases, such as horizontal and vertical scaling, database replication, and sharding, focusing on achieving fault tolerance and autoscaling capabilities. The authors emphasize the importance of making applications stateless by using external solutions for session storage and user files to facilitate horizontal scaling. The study concludes that to build a fault-tolerant web application, scalability is essential, and autoscaling can significantly reduce manual configuration time and improve responsiveness to traffic spikes. The research offers insights into how different scaling strategies can influence the performance and resilience of web applications, critical aspects also considered in the comparison between CSR and SSR.

\citet{clientServerSide:2020} investigated the performance of web applications built using CSR and SSR rendering techniques. The study developed a web application named Show Up using both approaches and compared their performance based on FCP, Speed Index, TTI, and FMP metrics. The results showed that SSR rendering outperformed CSR across all evaluated metrics, delivering a faster and smoother user experience. In addition, SSR scored higher on Google audit metrics for performance, accessibility, and best practices. The authors provide a discussion on different rendering strategies and their implications for application performance. This enables further investigation into other contributing factors, such as the impact of JSON and HTML data formats on user experience and web application efficiency, analyzed in the present study.
